\scanpage{45}
\mspage{5}
En effet, comme $G$ est plat sur $S$, on est ramené pour la première assertion au \deleted{critère différentiel d'ailleurs \illegible{1}}\inserted{cas $S=\operatorname{Spec}(k)$, et on peut supp. \unsure{même} $k$ alg. clos. Alors $G=G_{\mathrm{réd}}\times G_{\mathrm{inf}}$, où $G_{\mathrm{inf}}$ est un}
% slanted
\inserted{\unsure{p-groupe}, et $ {}_nG={}_nG_{\mathrm{réd}}$ \deleted{\unsure{où}} $n$ premier à $p$. On pourra \unsure{donc} supposer} \deleted{\unsure{le cas}} $G$ est lisse sur $k$, donc on peut \deleted{ramené au cas où} appliquer le critère différentiel : on est ramené à prouver que la \deleted{différentielle} application tangente est partout surjective, ce qui est immédiat, car elle s'identifie à la multiplication par $n$ [du fait que le foncteur $\operatorname{Lie}(G)$ est additif, en prenant les groupes alg. comm. sur $k$].

\uuline{Lemme} Soient $S$ schéma \deleted{de type} loc noeth. connexe, $X$ un \deleted{schéma $S$ \unsure{etc}} $S$-\deleted{\unsure{\illegible{1}|pré}}schéma \deleted{étale}\inserted{non ramifié}, $X_1$ et $X_2$ deux sous-schémas \deleted{fermés} de $X$ \uline{étales} \inserted{et \uline{finis}} sur $S$. \deleted{On suppose $X$, ou les $X_1$ et $X_2$, \uline{fini} sur $S$}. Soit $s\in S$ tel que $X_{1s}=X_{2s}$, alors $X_1=X_2$.

% slanted
\margin{Introduire $X_1\times_X X_2=X'$, fini et étale sur $S$, on est ramené au cas où $X_1\subset X_2$, à prouver que l'ens des pts où $X_{1s}=X_{2s}$ est ouvert et fermé \unsure{---}}

Supposons d'abord $S=\operatorname{Spec}(A)$, $A$ anneau local. On est ramené aussitôt au cas où $A$ est complet, donc \deleted{où $X$ est fini sur $S$}\inserted{[comme $X$ se décompose \deleted{de}\inserted{en $X_i$ qui sont} schémas locaux finis sur $S$]} \deleted{et schémas \unsure{éloignés} de la fibre pt sur $s$]} au cas où $X$ est fini sur $S$ [le cas où $X_i$ \deleted{\unsure{étale}} \unsure{au} vois. des pts sur $s$ étant \unsure{traité}]
% archive sheet number: 44
